\begin{lemma}[Second derivative of a penalised lift in a direction the penalty misses]
Let $P$ and $Q$ be real normed spaces, let $R$ be a real inner product space, let
$L:P\to Q$ and $M:P\to R$ be continuous linear maps, let $\lambda\in\R$, let $W\in P$,
and let $g:Q\to\R$ be $C^2$ near $L(W)$. Define
\[
  \Psi(U)=g\bigl(L(U)\bigr)+\lambda\,\frac{\|M(U)\|^2}{2}.
\]
Then for every $\zeta\in P$ with $M\zeta=0$,
\[
  D^2\Psi(W)[\zeta,\zeta]=D^2g\bigl(L(W)\bigr)\bigl[L\zeta,\ L\zeta\bigr].
\]
\end{lemma}
\begin{proof}
Write $q(u)=\|u\|^2/2$ for $u\in R$, so that the penalty is $\lambda\,(q\circ M)$. The
function $q$ is $C^2$ on all of $R$, because $u\mapsto\|u\|^2$ is a smooth function on an
inner product space; hence $q\circ M$ is $C^2$ on all of $P$, being the composition of a
$C^2$ function with a continuous linear map. Likewise $g\circ L$ is $C^2$ near $W$,
since $g$ is $C^2$ near $L(W)$ and $L$ is continuous linear.

\emph{Step 1: the first derivative splits on a neighbourhood of $W$.}
For every $U$ in a neighbourhood of $W$ both $g\circ L$ and $q\circ M$ are
differentiable at $U$; by the sum rule and the rule for a constant multiple,
\[
  D\Psi(U)=D(g\circ L)(U)+\lambda\,D(q\circ M)(U).
\]
Thus the two maps $U\mapsto D\Psi(U)$ and
$U\mapsto D(g\circ L)(U)+\lambda\,D(q\circ M)(U)$, both with values in the space of
continuous linear functionals on $P$, agree on a neighbourhood of $W$.

\emph{Step 2: differentiating once more.}
Because $g\circ L$ and $q\circ M$ are $C^2$ near $W$, their first derivatives are
differentiable at $W$, with derivatives $D^2(g\circ L)(W)$ and $D^2(q\circ M)(W)$. Two
maps that agree on a neighbourhood of $W$ have the same derivative at $W$, so by Step 1
\[
  D^2\Psi(W)=D^2(g\circ L)(W)+\lambda\,D^2(q\circ M)(W).
\]
Evaluating at the pair $(\zeta,\zeta)$,
\begin{equation}\label{eq:plsd-split}
  D^2\Psi(W)[\zeta,\zeta]
    =D^2(g\circ L)(W)[\zeta,\zeta]+\lambda\,D^2(q\circ M)(W)[\zeta,\zeta].
\end{equation}

\emph{Step 3: the penalty term vanishes.}
By the formula for the second derivative of a composition with a continuous linear map
on the right, applied to $q$ and $M$,
\[
  D^2(q\circ M)(W)[\zeta,\zeta]=D^2q\bigl(M(W)\bigr)\bigl[M\zeta,\ M\zeta\bigr].
\]
By hypothesis $M\zeta=0$, and $D^2q(M(W))$ is a bilinear map, hence vanishes as soon as
one of its arguments is $0$. Therefore $D^2(q\circ M)(W)[\zeta,\zeta]=0$.

\emph{Step 4: conclusion.}
Applying the same composition formula to $g$ and $L$ gives
$D^2(g\circ L)(W)[\zeta,\zeta]=D^2g(L(W))[L\zeta,L\zeta]$. Substituting this and
Step 3 into \eqref{eq:plsd-split} leaves
\[
  D^2\Psi(W)[\zeta,\zeta]=D^2g\bigl(L(W)\bigr)\bigl[L\zeta,\ L\zeta\bigr],
\]
as asserted.
\end{proof}
